\section{Quantization: choices and tests}\label{sec:quant}

\subsection{Which builds, and why}
The ladder is four Unsloth Dynamic builds between 16 and 24\,GiB (\cref{tab:arms}). The choice was
bounded from both sides. Below it, the 3-bit build was the fastest configuration measured in the whole
study and failed every agentic task it was given (\cref{sec:agentic}). Above it, an 8-bit build
(31\,GB) leaves no room for context and BF16 does not fit at all. \Qfour{} is Unsloth's recommended
default for this model; \QsixXL{} is the largest build that still leaves a usable window, which also
makes it the only available reference. Two further points, \IQfour{} and NVFP4, appear in the
historical rows. Six intermediate builds in the repository (\arm{Q4\_K\_S}, \arm{Q5\_K\_S},
\arm{Q5\_K\_M}, \arm{Q6\_K\_M}, \arm{Q6\_K\_L} and the deleted \arm{Q4\_K\_M}) were not tested.

\subsection{Divergence ranks the ladder}\label{sec:kld}

\begin{figure}[t]
\centering
\begin{tikzpicture}
\begin{axis}[paper,ybar,bar width=17pt,height=5.8cm,width=0.85\linewidth,
  xtick={1,2,3},
  xticklabels={WikiText-2 prose,Django code,HumanEval+ prompts},
  xmin=0.5,xmax=3.5,ymin=0,ymax=0.041,
  ylabel={mean KL divergence vs.\ \QsixXL},
  yticklabel style={/pgf/number format/fixed,/pgf/number format/precision=3},
  scaled y ticks=false,
  legend pos=north west,legend columns=1,legend image post style={scale=0.9},
  error bars/y dir=both,error bars/y explicit,error bars/error bar style={black,line width=0.6pt}]
\addplot[fill=cbsky,draw=cbblue!80!black] coordinates {(1,0.003321) +- (0,0.000126) (2,0.005829) +- (0,0.000233) (3,0.010403) +- (0,0.000448)};
\addplot[fill=cborange,draw=cborange!60!black,postaction={pattern=north east lines,pattern color=black!40}] coordinates {(1,0.004465) +- (0,0.000281) (2,0.010285) +- (0,0.000458) (3,0.017285) +- (0,0.000703)};
\addplot[fill=cbred,draw=cbred!70!black,postaction={pattern=crosshatch dots,pattern color=black!45}] coordinates {(1,0.008207) +- (0,0.000340) (2,0.021529) +- (0,0.000834) (3,0.036129) +- (0,0.001534)};
\addlegendimage{line legend,dashed,black!60,line width=0.6pt}
\legend{\Qsix,\Qfive,\Qfour,0.007 (external threshold)}
\draw[dashed,black!60] (axis cs:0.5,0.007) -- (axis cs:3.5,0.007);
\end{axis}
\end{tikzpicture}
\caption{Divergence separates every adjacent pair of builds in every domain, and the damage grows as
the corpus approaches the real task. Bars are mean KLD $\pm 1$ tool-reported SE; $n=65{,}536$ tokens
per cell (18{,}432 for task prompts); context 2{,}048; \flag{q4\_0} KV; greedy teacher forcing. The
dashed line is an external quality threshold, shown for orientation only.}\label{fig:kld}
\end{figure}

\Cref{fig:kld} is the central result. Against the \QsixXL{} reference, mean KLD on code is
$0.0058\pm0.0002$ for \Qsix, $0.0103\pm0.0005$ for \Qfive{} and $0.0215\pm0.0008$ for \Qfour. The order \Qsix{} $<$ \Qfive{} $<$ \Qfour{} holds in every domain (\cref{tab:paired-kld}). Each cell took about twenty minutes. The full table with quantiles is in
\cref{app:kld}.

Three things follow.

\textbf{Prose tables flatter the quantizer.} The same build diverges 1.8--2.6$\times$ more on code than
on prose, and 3.1--4.4$\times$ more on the HumanEval+ prompts than on prose. The amplification grows as
bits are removed (prose to task: $3.1\times$, $3.9\times$, $4.4\times$). Against the external 0.007
line, two of three builds pass on prose, one passes on code, none passes on the task prompts. Public
quantization tables are almost all measured on prose.

\textbf{Top-1 agreement points the other way.} On code every build agrees with the reference on the
argmax token \emph{more} often than on prose (98.4--99.1\,\% against 96.2--97.5\,\%), while its mean
KLD is double. Code is far more predictable (reference perplexity 1.18 against 5.79), so the top token
usually survives even as the distribution under it moves. A table reporting only top-1 agreement would
conclude this model quantizes better for code. Report the pair. On the task prompts \Qfour{} agrees on
95.9\,\% of tokens: about one token in 24 differs from the reference under greedy decoding.

\textbf{How sure is the ordering?} The tool's standard error treats 65{,}536 tokens from contiguous
windows as independent, and comparing two builds through their separate errors ignores that they were
scored on the same tokens against the same reference logits. Both problems have a fix that costs no
GPU time. The tool prints a running mean after each 2{,}048-token window, so per-window means can be
recovered by differencing, and the builds can be compared \emph{paired by window}
(\cref{tab:paired-kld}). Every adjacent difference has a bootstrap interval that excludes zero, and
the larger build has the lower divergence in 30 to 32 of 32 windows. The naive calculation had given
8.7--18 ``standard errors'' on code; the paired window-level $t$ is 4.1--4.5 there, and 5.5--11 on
prose. The ordering is secure. The size of the naive separations was not.

\begin{table}[t]
\centering\small
\caption{Paired differences in mean KLD between builds, by window. Per-window means recovered from the
tool's running mean (printed to five decimals); percentile bootstrap over windows, $B=20{,}000$. Windows
are contiguous chunks of one corpus per domain, so the intervals cover window-to-window variation on
these corpora, not the choice of corpus.}\label{tab:paired-kld}
\shrinktofit{\begin{tabular}{@{}llrlr@{}}
\toprule
corpus (windows) & pair & mean difference & bootstrap 95\,\% CI & windows where larger build is closer \\
\midrule
code (32) & \Qfive{} $-$ \Qsix & 0.00446 & [0.0027, 0.0065] & 32 / 32 \\
 & \Qfour{} $-$ \Qfive & 0.01124 & [0.0063, 0.0170] & 31 / 32 \\
prose (32) & \Qfive{} $-$ \Qsix & 0.00114 & [0.0008, 0.0016] & 30 / 32 \\
 & \Qfour{} $-$ \Qfive & 0.00375 & [0.0032, 0.0045] & 32 / 32 \\
task prompts (9) & \Qfive{} $-$ \Qsix & 0.00689 & [0.0056, 0.0081] & 9 / 9 \\
 & \Qfour{} $-$ \Qfive & 0.01884 & [0.0160, 0.0215] & 9 / 9 \\
\bottomrule
\end{tabular}}
\end{table}

\textbf{What this instrument measures.} Divergence here is a distance: the mean teacher-forced KL from
one 6-bit reference, on these corpora, at 2{,}048 tokens of context. It is not accuracy, and a larger
distance is not by itself a worse model; \citet{dutta2024accuracy} keep distance metrics and capability
metrics apart for the same reason. KL divergence obeys no triangle inequality, so an ordering of
distances from \QsixXL{} does not establish the ordering of distances from the unquantized model, and
errors that the reference shares with a build would make that build look closer than it is. The builds
also mix bit-widths per tensor, so the ordering describes three recipes, not a clean effect of bit
count. I use ``damage'' in this report as shorthand for this distance, and the task benchmarks below
are the only evidence here about whether it matters.

\subsection{The damage lives in a thin tail}\label{sec:tail}

\begin{figure}[t]
\centering
\begin{tikzpicture}
\begin{axis}[paper,height=5.8cm,ymode=log,
  symbolic x coords={median,p90,p95,p99,p99.9,max},xtick=data,
  xlabel={quantile of per-token KL divergence},
  ylabel={code $\div$ prose},
  ymin=0.003,ymax=20,legend pos=south east,
  extra y ticks={1},extra y tick labels={},extra y tick style={grid style={black!60,dashed}}]
\addplot[cbblue,mark=*] coordinates {(median,0.0050) (p90,0.449) (p95,1.527) (p99,4.988) (p99.9,6.352) (max,0.629)};
\addplot[cborange!85!black,mark=square*] coordinates {(median,0.0055) (p90,0.541) (p95,1.888) (p99,6.709) (p99.9,8.528) (max,0.813)};
\addplot[cbred,mark=triangle*] coordinates {(median,0.0048) (p90,0.613) (p95,2.184) (p99,8.090) (p99.9,8.448) (max,0.730)};
\legend{\Qsix,\Qfive,\Qfour}
\node[font=\scriptsize,text=black!70,anchor=west] at (axis cs:median,1.6) {code worse $\uparrow$};
\node[font=\scriptsize,text=black!70,anchor=west] at (axis cs:median,0.6) {prose worse $\downarrow$};
\end{axis}
\end{tikzpicture}
\caption{``Code is twice as bad as prose'' is an average of two opposite facts. At the median a code
token is perturbed about $200\times$ \emph{less} than a prose token; the order flips between the 90th
and 95th percentile; at the 99th, code is 5--8$\times$ worse. The single worst token is again milder on
code. Quantiles carry no interval, and the median row cannot rank the builds (its inputs have one or
two significant digits).}\label{fig:tail}
\end{figure}

\Cref{fig:tail} divides each quantile of per-token KLD on code by the same quantile on prose. About
90\,\% of code tokens are nearly untouched by quantization; a band between the 95th and 99.9th
percentile moves a great deal. A mean-only report averages ``almost nothing happens'' with ``something
drastic happens'' and states neither.

This is a candidate explanation for why benchmarks see so little. A benchmark scores outcomes, not
tokens. A perturbation confined to a few percent of positions changes an outcome only when one of
those positions is decisive. That predicts small discordance between builds, and small discordance is
what the paired benchmarks below show (3 and 5 problems of 164). It is a hypothesis with one successful
quantitative prediction; no experiment here manipulates the tail and watches a benchmark move.

\paragraph{A scope I got wrong at first.} I initially read the tail as something quantization
\emph{does}. The study contains a control that says otherwise. One cell perturbs only the KV cache
(\flag{q4\_0} against \flag{f16}) on the reference weights. Normalising each cell by its own
mean, the 99th percentile sits at 24--26$\times$ the mean for all three quantized builds on code, at
23$\times$ for the KV-only cell, and at 8.4--8.6$\times$ for every build on prose
(\cref{tab:tailshape}). The \emph{shape} of the tail belongs to the corpus. Quantization sets the
\emph{level}: $0.0030 \rightarrow 0.0058 \rightarrow 0.0103 \rightarrow 0.0215$. The control was not
designed for this purpose, so the comparison is suggestive, not clean.

\begin{table}[t]
\centering\small
\caption{Tail shape, each cell normalised by its own mean KLD. Flat within a corpus over a $7\times$
range of mean divergence; different by about $3\times$ between corpora; reproduced by a perturbation
with no weight quantization in it.}\label{tab:tailshape}
\shrinktofit{\begin{tabular}{@{}lllrrrr@{}}
\toprule
cell & corpus & perturbation & mean KLD & median/mean & p95/mean & \textbf{p99/mean} \\
\midrule
\Qsix / \Qfive / \Qfour & prose & weights & .0033 / .0045 / .0082 & 0.41--0.43 & 3.1--3.2 & \textbf{8.4--8.6} \\
\Qsix / \Qfive / \Qfour & code & weights & .0058 / .0103 / .0215 & 0.001 & 2.6--2.8 & \textbf{24.4--26.2} \\
\QsixXL & code & \textbf{KV dtype only} & .0030 & 0.001 & 2.8 & \textbf{22.9} \\
\Qsix / \Qfive / \Qfour & task prompts & weights & .0104 / .0173 / .0361 & 0.02 & 4.8--4.9 & \textbf{17.1--18.3} \\
\bottomrule
\end{tabular}}
\end{table}

\subsection{Perplexity: right order, no resolution, and a 36-fold swing}\label{sec:ppl}

\begin{table}[t]
\centering\small
\caption{WikiText-2 perplexity\hist, Protocol~1: 602 windows of 512 tokens, scoring the second half of
each window with the first half as context; $\pm$ is the tool's standard error. The whole ladder spans
0.033 while each point carries $\pm0.041$.}\label{tab:ppl}
\shrinktofit{\begin{tabular}{@{}lrrr@{}}
\toprule
build & PPL & $\pm$ SE & $\Delta$ vs best \\
\midrule
\QsixXL & 6.6511 & 0.0411 & --- \\
\Qfive & 6.6556 & 0.0411 & +0.07\,\% \\
\Qfour & 6.6617 & 0.0412 & +0.16\,\% \\
\IQfour & 6.6839 & 0.0413 & +0.49\,\% \\
\emph{NVFP4 (vLLM), re-scored on the same windows} & 6.7073 & n/r & +0.85\,\% \\
\bottomrule
\end{tabular}}
\end{table}

Perplexity orders the ladder correctly (\cref{tab:ppl}) and cannot certify a single step of it: even
the extremes overlap. A second protocol from the same period (20 windows of 4{,}096 tokens, f16 KV)
gave different absolute values (5.50--5.55) and was not even monotone. The KV-cache experiment makes
the same point from the other side: switching the cache from \flag{f16} to \flag{q4\_0} moved code
perplexity from 1.1791 to 1.1809 (+0.15\,\%) while the measured KL divergence of that change is half
that of dropping a whole quantization level.

\begin{figure}[t]
\centering
\begin{tikzpicture}
\begin{axis}[paper,xbar,height=4.4cm,width=0.62\linewidth,bar width=11pt,
  ytick={1,2,3},ymin=0.4,ymax=3.6,
  yticklabels={{matched: same file, 602 windows, half-window scoring},{same file and windows, all positions scored},{as first measured: other file, 160 windows, no context}},
  yticklabel style={font=\scriptsize,text width=5.4cm,align=right},
  xmin=0,xmax=42,xlabel={apparent degradation vs.\ \QsixXL{} (\%)},
  nodes near coords,nodes near coords style={font=\scriptsize,anchor=west},
  point meta=explicit symbolic]
\addplot[fill=cbred!80,draw=cbred!60!black] coordinates {(0.85,1)[{PPL 6.71 (+0.8\,\%)}] (21.4,2)[{PPL 8.08 (+21\,\%)}] (29.1,3)[{PPL 8.58 (+29\,\%)}]};
\end{axis}
\end{tikzpicture}
\caption{One checkpoint, one benchmark corpus, three conventions\hist. The NVFP4 build looked 29\,\%
worse than the GGUF ladder until it was re-scored on the ladder's exact windows with the ladder's
scoring rule. The corpus file (two conversions of the same dataset, 297K against 309K tokens), the
number of windows and the scoring rule all contributed; tokenizer mismatch was suspected, tested and
ruled out.}\label{fig:swing}
\end{figure}

\Cref{fig:swing} is the most transferable result of the exploratory period. The first NVFP4 perplexity
(8.58 against 6.65) suggested a 29\,\% loss. Re-scored on the same windows under the same half-window
rule it is 6.71: a 0.8\,\% loss, a 36-fold change in the estimated effect from convention alone. The
corrected number still ranks NVFP4 below every GGUF build, so the ranking held and the magnitude did
not. A perplexity published without its corpus file, window count and scoring rule cannot be compared
with another.

\subsection{What task benchmarks can say: bounds}\label{sec:bounds}

\begin{table}[t]
\centering\small
\caption{HellaSwag, first battery. 400 items, the same items for every build, greedy log-likelihood
scoring. The most quantized build scores nominally highest; two builds answer all 400 items
identically.}\label{tab:hellaswag}
\shrinktofit{\begin{tabular}{@{}lrrl@{\hskip 1.5em}lrrr@{}}
\toprule
build & correct & acc.\,\% & Wilson 95\,\% & pair & discordant & diff.\ (pts) & min.\ attainable $p$ \\
\midrule
\QsixXL & 331 & 82.75 & [78.7, 86.1] & \QsixXL{} vs \Qsix & 2 & $+0.50$ & 0.50 \\
\Qsix & 329 & 82.25 & [78.2, 85.7] & \QsixXL{} vs \Qfive & \textbf{0} & 0.00 & 1.00 \\
\Qfive & 331 & 82.75 & [78.7, 86.1] & \QsixXL{} vs \Qfour & 2 & $-0.50$ & 0.50 \\
\Qfour & \textbf{333} & \textbf{83.25} & [79.3, 86.6] & \Qsix{} vs \Qfour & 4 & $-1.00$ & 0.125 \\
\bottomrule
\end{tabular}}
\end{table}

\begin{table}[t]
\centering\small
\caption{HumanEval+, first battery, paired. 164 problems, Qwen's official non-thinking sampling, same
seed, context 32{,}768, \textbf{no speculation on either arm}, zero empty responses.}\label{tab:he-paired}
\shrinktofit{\begin{tabular}{@{}lrlrl@{}}
\toprule
build & HumanEval & Wilson 95\,\% & HumanEval+ & Wilson 95\,\% \\
\midrule
\QsixXL & 94.51\,\% (155) & [89.9, 97.1] & 90.24\,\% (148) & [84.7, 93.9] \\
\Qfour  & 93.90\,\% (154) & [89.1, 96.7] & 89.63\,\% (147) & [84.0, 93.4] \\
\midrule
\multicolumn{5}{@{}l}{paired, \Qfour{} $-$ \QsixXL: base $-0.61$ pts, 95\,\% score CI $[-3.80,+2.28]$, 3 discordant (1 vs 2), min.\ $p$ 0.25} \\
\multicolumn{5}{@{}l}{\phantom{paired, \Qfour{} $-$ \QsixXL: }plus $-0.61$ pts, 95\,\% score CI $[-4.19,+2.74]$, 5 discordant (2 vs 3), min.\ $p$ 0.0625} \\
\bottomrule
\end{tabular}}
\end{table}

\Cref{tab:hellaswag,tab:he-paired} show the two first-battery task benchmarks on the same builds that
divergence orders in every window. HellaSwag spreads the four builds over one point, puts
the 4-bit build on top, and finds two builds identical on every item. HumanEval+ has the ladder's
extremes agreeing on 159 of 164 problems.

My first write-up reported the HumanEval+ result as ``McNemar $p=1.0$, indistinguishable''. That test
could not have reached significance: with three discordant pairs the smallest possible two-sided $p$
is 0.25. This is the very error the report warns about, and I made it in the report's own anchor
experiment. The informative statement is the interval:

\begin{takeaway}
A $3.7\times$ increase in code-prompt KL divergence (\QsixXL{} $\rightarrow$ \Qfour) moves HumanEval+
pass@1 by $-0.6$ points, with a 95\,\% confidence interval reaching about 4 points. Benchmarks at these sample
sizes \emph{bound} quantization damage; they do not \emph{rank} it.
\end{takeaway}

\begin{table}[t]
\centering\small
\caption{HumanEval / HumanEval+ pass@1\hist, 164 problems, greedy. Monotone in bit-width, but only the
3-bit step exceeds the $\pm4.6$-point interval. With reasoning on (4{,}096-token budget) every build
scores lower, and the builds differ mainly in how often they run out of budget and return nothing.
MTP and DFlash2 rows are not comparable at a fixed build (\cref{sec:nonequiv}).}\label{tab:humaneval}
\shrinktofit{\begin{tabular}{@{}llrrlr@{}}
\toprule
mode & build (speculation) & HE & HE+ & Wilson 95\,\% (HE+) & empty \\
\midrule
non-thinking & \Qthree & 84.1 & 81.7 & [75.1, 86.9] & 0\,\% \\
 & \IQfour & 90.2 & 87.8 & [81.9, 92.0] & 0\,\% \\
 & \Qfive & 93.3 & 90.9 & [85.5, 94.4] & 0\,\% \\
 & \QsixXL & 93.9 & 91.5 & [86.2, 94.9] & 0\,\% \\
 & \QsixXL, window 32K \emph{and} 131K (control) & 93.3 & 90.2 & [84.7, 93.9] & 0\,\% \\
\addlinespace
thinking & \IQfour{} (MTP) & 86.6 & 86.0 & [79.8, 90.5] & 12.8\,\% \\
 & \Qfour{} (MTP) & 87.8 & 86.0 & [79.8, 90.5] & 12.2\,\% \\
 & \Qfive{} (MTP) & 89.0 & 86.0 & [79.8, 90.5] & 11.0\,\% \\
 & \QsixXL{} (MTP) & 90.9 & 88.4 & [82.6, 92.5] & 7.9\,\% \\
 & \IQfour{} (DFlash2) & 89.0 & 87.8 & [81.9, 92.0] & 10.4\,\% \\
 & \Qfour{} (DFlash2) & 90.2 & 87.2 & [81.2, 91.5] & 8.5\,\% \\
 & NVFP4 on vLLM (none) & 85.4 & 84.1 & [77.8, 89.0] & 12.8\,\% \\
\bottomrule
\end{tabular}}
\end{table}

The historical HumanEval+ ladder (\cref{tab:humaneval}) agrees. Scores rise with bit-width in the same
order as perplexity and divergence, yet the three upper builds span 3.7 points against a $\pm4.6$-point
interval; the 5-bit to 6-bit step is one problem. Two side results are worth keeping. The context
control is the cleanest null in the study: allocating a 131K window instead of 32K changed nothing, to
the last problem. And in thinking mode with a 4{,}096-token budget, three of four MTP rows score exactly
86.0 on HumanEval+ and differ only in their empty-response rate (12.8\,\% down to 7.9\,\% with
fidelity): that benchmark partly measures whether reasoning fits its budget. The same mechanism
returns in \cref{sec:ruler}.

\subsection{The 4-bit KV cache is validated at 2K tokens, and only there}\label{sec:kvq}
Every context result in this report depends on a quantized KV cache, so its cost was measured:
\flag{q4\_0} against \flag{f16} on the reference build, code corpus, gives mean KLD
$0.00296\pm0.00013$ with 99.4\,\% top-1 agreement. That is about half the size of the step from
\QsixXL{} to \Qsix{} (0.0058). It is a comparison of sizes; KL divergences from a common reference do
not add, and no cell applies both perturbations.

This number should not be carried to long context. It was taken at 2{,}048 tokens, where the cache
barely matters. The literature says 4-bit KV needs per-channel key quantization to stay near-lossless
\citep{liu2024kivi,hooper2024kvquant}, which \flag{q4\_0} lacks, and a BF16-referenced community
measurement on the sibling Qwen3.6-27B over about 250K tokens reports KLD of 0.024 for \flag{q8\_0} and
0.09--0.12 for \flag{q4\_0} \citep{localbench2026kvquant}: one to two orders of magnitude above my
short-context figure. I could not measure divergence at depth on this host (\cref{sec:hardware}). The
second battery handles the question differently: it finds the gentlest KV precision each build can
hold at the full window (\cref{sec:kvmap}) and measures task accuracy at that precision
(\cref{sec:aalcr}).
